\documentclass[exercice,a4]{thedocument}%
\usepackage[french]{babel}%
\usepackage{tabularx}%
\usepackage{array}%
\startdatakeys
    \source{22GENMATMEAG2}
    \BO{2026}
    \cycle{4}
    \competencesDuSocle{CA,RA}
    \objectifsApprentissage{C4C2aa,C4C2ab,C4D2ae,C4C2ag}
    \motsCles{médiane,tableur,moyenne,pourcentage}
\enddatakeys
\begin{document}%

\medskip

L'artiste français Jean Lurçat a produit dix tapisseries de surfaces différentes, exposées dans la ville d'Angers.

La surface approximative de chacune de ces tapisseries a été saisie dans la feuille de calcul ci-dessous. (source: https://musees.angers.fr)

\begin{center}
\begin{tabularx}{0.6\linewidth}{|c|c|>{\centering \arraybackslash}X|}\hline
1&\textbf{Nom de la tapisserie}			&\textbf{Surface (en m$^2$)}\\ \hline
2&La Grande Menace				&39,2\\ \hline
3&L'homme d'Hiroshima			&12,8\\ \hline
4&Le Grand charnier				&32,4\\ \hline
5&La Fin de tout				&10,2\\ \hline
6&L'Homme en gloire dans la paix&57,5\\ \hline
7&L'eau et le feu				&27,2\\ \hline
8&Champagne						&30,9\\ \hline
9&La Conquête de l'espace		&45,9\\ \hline
10&La Poésie					&45,4\\ \hline
11&Ornamentos sagrados			&45,4\\ \hline
12&\textbf{Total}				&\\ \hline
\end{tabularx}
\end{center}

\begin{enumerate}
\item Quelle tapisserie a la surface maximale ? On ne demande pas de justifier.
\item Vérifier que la surface médiane de ces tapisseries est égale à 35,8~m$^2$.
\item Quelle formule peut-on saisir dans la cellule B12 afin de calculer la surface totale de ces
dix tapisseries ?

On ne demande pas de justifier.
\item Calculer la surface moyenne de ces tapisseries.
\item Sara affirme: \og 40\% de ces tapisseries ont une surface supérieure à 45 m$^2$ \fg.

Sara a-t-elle raison ? Justifier.
\end{enumerate}

\bigskip

\showthesource
\end{document}
